% Part: first-order-logic
% Chapter: completeness
% Section: introduction

\documentclass[../../../include/open-logic-section]{subfiles}

\begin{document}

\iftag{FOL}
      {\olfileid{fol}{com}{int}}
      {\olfileid{pl}{com}{int}}

\olsection{مقدمة}

مبرهنة الاكتمال من أهم النتائج الأساسية في المنطق. ولها صيغتان سنبرهن
تكافؤهما. تقول الصيغة الأولى أمرًا أساسيًا عن العلاقة بين النتيجة
الدلالية ونسق ال!!{derivation} لدينا: إذا لزمت !!a{sentence}~$!A$ عن
بعض ال!!{sentence}s $\Gamma$، فهناك أيضًا !!a{derivation} يثبت
$\Gamma \Proves !A$. وهكذا يكون نسق ال!!{derivation} قويًا إلى أقصى
حد ممكن من غير أن يثبت أشياء لا تلزم فعلًا.

أما الصيغة الثانية فيمكن عرضها كنتيجة عن وجود النماذج: كل مجموعة متسقة
من ال!!{sentence}s قابلة للإرضاء. والاتساق مفهوم نظري برهاني؛ إذ يقول
إن نسق ال!!{derivation} لدينا عاجز عن إنتاج !!{derivation}s من أنواع
معينة. ولكن مَن يضمن أن مجرد عدم وجود !!{derivation}s من نوع معين من
~$\Gamma$ يعني وجود
\iftag{FOL}{!!a{structure}~$\Struct{M}$}{!!a{valuation}~$\pAssign{v}$}
تحقق $\iftag{FOL}{\Sat{M}}{\pSat{v}}{\Gamma}$؟ قبل أن تُبرهن مبرهنة
الاكتمال لأول مرة---بل قبل أن تتوافر لنا أنساق ال!!{derivation} التي
لدينا الآن---كان عالم الرياضيات الألماني الكبير ديفيد هلبرت يرى أن
اتساق النظريات الرياضية يضمن وجود الأشياء التي تتحدث عنها. وقد عبّر عن
ذلك في رسالة إلى غوتلوب فريغه بقوله:
\begin{quote}
  إذا كانت البديهيات المعطاة اعتباطًا، مع كل نتائجها، لا يناقض بعضها
  بعضًا، فهي صادقة، والأشياء التي تعرّفها البديهيات موجودة. وهذا هو عندي
  معيار الصدق والوجود.
\end{quote}
وخالفه فريغه بشدة. وتبين الصيغة الثانية لمبرهنة الاكتمال أن هلبرت كان
محقًا، على الأقل بالمعنى القائل إنه إذا كانت البديهيات متسقة، وُجدت
\emph{بعض} ال\iftag{FOL}{!!{structure}}{!!{valuation}} التي تجعلها
كلها صادقة.

وليست هذه الأسباب الوحيدة لأهمية مبرهنة الاكتمال---أو بالأحرى برهانها.
فلها نتائج مهمة عديدة، سنناقش بعضها على حدة. فمثلًا، لما كان كل
!!{derivation} يبين أن $\Gamma \Proves !A$ منتهيًا، ولا يستطيع لذلك أن
يستعمل إلا عددًا منتهيًا من ال!!{sentence}s في~$\Gamma$، يترتب من
مبرهنة الاكتمال أنه إذا كانت $!A$ نتيجة لـ~$\Gamma$، فهي نتيجة بالفعل
لمجموعة جزئية منتهية من~$\Gamma$. ويسمى هذا \emph{التراص}. وبصورة
مكافئة، إذا كانت كل مجموعة جزئية منتهية من $\Gamma$ متسقة، وجب أن تكون
$\Gamma$ نفسها متسقة.

ومع أن مبرهنة التراص تلزم عن مبرهنة الاكتمال بطريق يمر عبر
ال!!{derivation}s، فمن الممكن أيضًا استعمال \emph{برهان} مبرهنة
الاكتمال لإثباتها مباشرة. فالبرهان يأخذ مجموعة من ال!!{sentence}s لها
خاصية معينة---هي الاتساق---وينشئ من هذه المجموعة
\iftag{FOL}{!!a{structure}}{!!a{valuation}} لها خصائص معينة (وفي هذه
الحالة أنها ترضي المجموعة). ويمكن استعمال البناء نفسه تقريبًا لإثبات
التراص مباشرة، بالبدء بمجموعات من ال!!{sentence}s «القابلة للإرضاء
منتهيًا» بدلًا من المجموعات المتسقة.
\iftag{FOL}{ويثمر البناء أيضًا نتائج أخرى؛ منها مثلًا أن لكل مجموعة
  قابلة للإرضاء من ال!!{sentence}s نموذجًا منتهيًا أو !!{denumerable}.
  (وتسمى هذه النتيجة مبرهنة لوفنهايم--سكولم.) وبوجه عام، يكثر في المنطق
  استعمال إنشاء ال!!{structure}s من مجموعات ال!!{sentence}s، بل يستعمل
  أحيانًا في الفلسفة أيضًا.}{}

\end{document}
